\subsection{NanoMem Training Procedure}
\label{app:nanomem-training-procedure}

This subsection summarizes the online GRPO training loop used to optimize only
the Synthesizer/Compressor policy in \sys. Planner, preprocessor, retriever,
answerer, judge, and reference policy are frozen throughout training.

\begin{tcolorbox}[
  enhanced jigsaw,
  breakable,
  colback=white,
  colframe=black,
  boxrule=0pt,
  borderline north={0.7pt}{0pt}{black},
  borderline south={0.7pt}{0pt}{black},
  sharp corners,
  left=0pt,
  right=0pt,
  top=2pt,
  bottom=2pt,
  before skip=0.6em,
  after skip=0.8em,
]
\scriptsize
\noindent\textbf{Algorithm 1} GRPO Training for Multi-turn Memory Evidence Synthesis

\vspace{0.15em}
\hrule
\vspace{0.45em}

\noindent\textbf{Input:}\\
\hspace*{1.5em}Training set $\mathcal{D}_{\mathrm{train}}$; each sample contains query $q$, gold answer $a^\star$,\\
\hspace*{1.5em}gold evidence sessions $\mathcal{E}^\star$, source DB path, user id, and analysis metadata\\
\hspace*{1.5em}Trainable compressor policy $\pi_\theta$ initialized from Qwen3.5-9B\\
\hspace*{1.5em}Frozen planner, retriever, answerer, judge, and reference policy $\pi_{\mathrm{ref}}$\\
\hspace*{1.5em}Rollout group size $G=8$, maximum search rounds $R_{\max}=3$

\vspace{0.25em}
\noindent\textbf{Ensure:}\\
\hspace*{1.5em}Optimized compressor policy $\pi_{\theta'}$

\vspace{0.35em}
\begin{tabular}{@{}r@{\quad}p{0.84\linewidth}@{}}
1: & \textbf{for each} training step \textbf{do} \\
2: & \hspace*{1em}Sample one training example $x=(q,a^\star,\mathcal{E}^\star,\mathrm{DB},u,\mathrm{meta})$ from $\mathcal{D}_{\mathrm{train}}$ \\
3: & \hspace*{1em}\textbf{Stage 1: Input and initialization} \\
4: & \hspace*{2em}Load the sample-specific memory DB and user id; keep planner, retriever, answerer, judge, and $\pi_{\mathrm{ref}}$ frozen \\
5: & \hspace*{1em}\textbf{Stage 2--3: Online multi-turn group rollout} \\
6: & \hspace*{2em}\textbf{for} rollout $g=1,\ldots,G$ \textbf{do} \\
7: & \hspace*{3em}Initialize previous events $\mathcal{M}^{(0)}\leftarrow\emptyset$, previous cues $\mathcal{C}^{(0)}\leftarrow\emptyset$, and terminal flag $\mathrm{stop}\leftarrow\mathrm{false}$ \\
8: & \hspace*{3em}\textbf{for} round $r=1,\ldots,R_{\max}$ \textbf{do} \\
9: & \hspace*{4em}Use the frozen planner to generate retrieval cues $\mathcal{C}^{(r)}$ from $q$, $\mathcal{C}^{(r-1)}$, and compressor reasoning when available \\
10: & \hspace*{4em}Use the frozen retriever to search sessions $\mathcal{S}^{(r)}$ from the sample-specific memory DB \\
11: & \hspace*{4em}Construct temporal hints online from the query and retrieved sessions \\
12: & \hspace*{4em}Sample compressor XML $z_g^{(r)}\sim\pi_\theta(\cdot\mid q,\mathrm{hints},\mathcal{M}^{(r-1)},\mathcal{S}^{(r)})$ \\
13: & \hspace*{4em}Parse $z_g^{(r)}$ into reasoning, events $\mathcal{M}^{(r)}$, and verdict $v^{(r)}$ \\
14: & \hspace*{4em}\textbf{if} $v^{(r)}=\mathrm{sufficient}$, no new evidence is found, or XML is invalid \textbf{then} set $\mathrm{stop}\leftarrow\mathrm{true}$ and break \\
15: & \hspace*{3em}\textbf{end for} \\
16: & \hspace*{3em}Use the frozen answerer to generate final answer $\hat{a}_g$ from $q$ and terminal events $\mathcal{M}^{(r)}$ \\
17: & \hspace*{3em}Store rollout trace $\tau_g$: retrieved sessions, compressor XML, reasoning, events, verdicts, terminal reason, and answer \\
18: & \hspace*{2em}\textbf{end for} \\
19: & \hspace*{1em}\textbf{Stage 4: Reward computation} \\
20: & \hspace*{2em}\textbf{for} rollout $g=1,\ldots,G$ \textbf{do} \\
21: & \hspace*{3em}Compute format reward from XML validity and schema compliance \\
22: & \hspace*{3em}Compute verdict reward by comparing sufficient/insufficient decisions with gold evidence coverage $\mathcal{E}^\star$ \\
23: & \hspace*{3em}Use the frozen judge to compare $\hat{a}_g$ with $a^\star$ and assign outcome reward \\
24: & \hspace*{3em}Compute length reward from compressor output compactness, mainly active when the outcome is correct \\
25: & \hspace*{3em}Combine format, verdict, outcome, and length rewards into scalar reward $R_g$ \\
26: & \hspace*{2em}\textbf{end for} \\
27: & \hspace*{1em}\textbf{Stage 5: GRPO policy update} \\
28: & \hspace*{2em}Compute group-relative advantages from $\{R_g\}_{g=1}^{G}$ for the same query \\
29: & \hspace*{2em}Increase likelihood of higher-reward compressor outputs and decrease likelihood of lower-reward outputs \\
30: & \hspace*{2em}Apply KL regularization against frozen reference policy $\pi_{\mathrm{ref}}$ \\
31: & \hspace*{2em}Update only compressor policy parameters $\theta$ \\
32: & \textbf{end for} \\
33: & Select the final checkpoint using training reward only; held-out sets are used for reporting \\
34: & \textbf{return} optimized compressor policy $\pi_{\theta'}$
\end{tabular}
\end{tcolorbox}

\subsection{Reward Design}
\label{app:reward-design}

We define the reward over a complete rollout trajectory
$\tau=\{(h_1,z_1),\ldots,(h_T,z_T)\}$, where $h_t$ is the retrieval context and
$z_t=(r_t,\mathcal{E}_t,\hat{v}_t)$ is the Synthesizer action at round $t$. For
coefficient analysis, the same trajectory is abstracted as
$\tau=(f,T,\mathbf{v},o,\ell)$, where $f$ is format validity, $T$ is the number
of retrieval rounds, $\mathbf{v}$ is the sequence of verdict correctness
signals, $o$ is final answer correctness, and $\ell$ is evidence compactness.

\noindent\textbf{Format reward.}
Let $f(\tau)\in\{0,1\}$ indicate whether all Synthesizer outputs in the rollout
are parseable under the required XML schema. Format validity acts as a gate:
if $f(\tau)=0$, all other reward terms are not evaluated and the trajectory
receives the fixed reward $c<0$. Equivalently, the format-failure penalty
magnitude is $-c>0$.

\noindent\textbf{Length reward.}
Let $m(z_t)$ be the token length of the Synthesizer output at round $t$.
The compactness score $\ell(\tau)\in[0,1]$ is activated only when the final
answer is correct:
\[
\ell(\tau)
=
\mathbb{I}\{R_{\mathrm{out}}(\tau)=1\}
\cdot
\mathrm{clip}\!\left(
1-\frac{1}{Tm_{\max}}\sum_{t=1}^{T}m(z_t),
0,1
\right),
\]
where $m_{\max}$ is the length threshold. This term prefers shorter evidence
among trajectories that already solve the task.

\noindent\textbf{Verdict reward.}
The Synthesizer must decide whether the current retrieved evidence covers the
query or whether a missing evidence gap remains. Let $\mathcal{G}_q$ denote the
gold evidence session set for query $q$, and let
$\mathcal{S}^{\mathrm{ret}}_{\le t}$ denote the session identifiers covered by
the accumulated retrieved sessions $\widetilde{\mathcal{S}}_{\le t}$. The
round-level verdict reward is
\[
v_t=
\begin{cases}
+1, & \mathcal{G}_q\subseteq\mathcal{S}^{\mathrm{ret}}_{\le t}
\ \text{and}\ \hat{v}_t=\texttt{sufficient},\\
+1, & \mathcal{G}_q\not\subseteq\mathcal{S}^{\mathrm{ret}}_{\le t}
\ \text{and}\ \hat{v}_t=\texttt{insufficient},\\
-1, & \text{otherwise}.
\end{cases}
\]
This term evaluates whether the policy correctly decides to stop or continue
under the objective retrieval state. Since it is computed independently at each
round, the trajectory-level verdict score is $V=\sum_{t=1}^{T}v_t$, providing
fine-grained process supervision for multi-round search.

\noindent\textbf{Outcome reward.}
The outcome reward ties the final Temporal Evidence Pool to task success. Let
$H_{T+1}$ be the Temporal Evidence Pool after the final round. The frozen answer
model generates $\hat{a}=\mathrm{Ans}(q,H_{T+1})$, and the evaluator checks
semantic equivalence with the gold answer $a^\star$:
\[
o(\tau)=\mathbb{I}\{\mathrm{Eval}(\hat{a},a^\star)=1\}.
\]
This component measures whether the constructed evidence state is sufficient
for answering the original query.

Combining these components, the total reward is
\[
R(\tau)
=
\begin{cases}
c, & f(\tau)=0,\\
w_o o(\tau)+w_v\sum_{t=1}^{T}v_t+w_l\ell(\tau), & f(\tau)=1,
\end{cases}
\]
where $w_o,w_v,w_l>0$ are the weights for outcome, verdict, and compactness.
Thus, $o$ and $\ell$ are trajectory-level signals, while $v_t$ is accumulated
over retrieval rounds and format validity gates the whole trajectory. We do not
introduce a separate event-time reward: incorrect event-time reasoning is
reflected in $o(\tau)$ through retrieval drift and final answer failure.

\subsection{Reward Coefficient Analysis}
\label{app:reward-design-coefficients}

\noindent\textbf{Problem Definition.}
We abstract a rollout trajectory as
$\tau=(f,T,\mathbf{v},o,\ell)$, where $f\in\{0,1\}$ denotes whether the
Synthesizer output is parseable under the required XML schema,
$T\in\{1,\ldots,T_{\max}\}$ is the number of retrieval rounds, and
$\mathbf{v}=(v_1,\ldots,v_T)$ records round-level verdict correctness with
$v_t\in\{-1,+1\}$. The final answer correctness is $o\in\{0,1\}$, and
$\ell\in[0,1]$ measures evidence compactness, with $\ell=0$ when $o=0$ because
compactness is only rewarded for correct outcomes. Let
$V=\sum_{t=1}^{T}v_t$ be the cumulative verdict score. Format acts as a gate:
if $f=0$, verdict, outcome, and length terms are not evaluated and the rollout
receives a fixed penalty $c$. Otherwise, the reward is:
\[
R(\tau)=
\begin{cases}
c, & f=0,\\
w_vV+w_o o+w_l\ell, & f=1,
\end{cases}
\]
where $c\in\mathbb{R}$ is the format-failure penalty and
$w_v,w_o,w_l>0$ are reward weights for verdict correctness, answer correctness,
and compactness, respectively. With $T_{\max}=3$, the wrong-outcome reward range
under format success is $[-3w_v,3w_v]$, while the correct-outcome reward range
is $[-3w_v+w_o,3w_v+w_o+w_l]$.

We require this reward to satisfy six ordering goals. (G1) \textit{Format
gate}: every format-valid trajectory should exceed format failure. (G2)
\textit{Outcome dominance}: any correct-outcome trajectory should exceed any
wrong-outcome trajectory, regardless of verdict history or compactness. (G3)
\textit{Verdict discrimination}: for fixed outcome and compactness, a larger
cumulative verdict score should receive larger reward. (G4) \textit{Multi-round
credit}: a trajectory with more correct verdict rounds should receive more
verdict credit. (G5) \textit{Compactness preference}: for fixed outcome and
verdict, greater compactness should receive larger reward. (G6)
\textit{Verdict over length}: a single verdict correction should dominate the
full compactness range, so process supervision cannot be overridden by the
length term.

We also define four guaranteed margins to express the intended priority among
learning signals. $\Delta_1$ is the \emph{format-gate margin}: it asks how far
the worst format-success trajectory remains above format failure. Format
failure receives the fixed reward $c$, while the worst parseable trajectory has
all other signals minimized, namely $V=-3$, $o=0$, and $\ell=0$:
\[
R^{\inf}_{f=1}=w_v(-3)+w_o\cdot 0+w_l\cdot 0=-3w_v,
\]
\[
\Delta_1=R^{\inf}_{f=1}-c=-3w_v-c.
\]
Thus, $\Delta_1>0$ guarantees that any parseable output is preferred to a
format failure. $\Delta_2$ is the \emph{outcome-discrimination margin}: under
format success, it compares the worst correct-answer trajectory with the best
wrong-answer trajectory. The former has $V=-3$ and $\ell=0$, while the latter
has $V=3$:
\[
R^{\inf}_{o=1}=-3w_v+w_o,\quad
R^{\sup}_{o=0}=3w_v,
\]
\[
\Delta_2=R^{\inf}_{o=1}-R^{\sup}_{o=0}=w_o-6w_v.
\]
The term $6w_v$ is the maximal swing introduced by cumulative verdict scoring;
$\Delta_2>0$ means that answering correctly remains better than answering
incorrectly even under this extreme verdict contrast. $\Delta_3$ is the
\emph{single-step verdict signal}: changing one round-level verdict from
incorrect to correct changes $v_t$ from $-1$ to $+1$, increasing $V$ by $2$:
\[
w_v(V+2)-w_vV=2w_v,
\]
\[
\Delta_3=2w_v.
\]
This margin is independent of $T$, $o$, and $\ell$. Finally, $\Delta_4$ is the
\emph{compactness range}: since $\ell\in[0,1]$ and the length term is
$w_l\ell$, the maximum reward variation induced by compactness is
\[
w_l\cdot 1-w_l\cdot 0=w_l,
\]
\[
\Delta_4=w_l.
\]
We require $\Delta_1>\Delta_2>\Delta_3>\Delta_4>0$ so that the guaranteed
discrimination decreases along the intended priority order: format validity,
answer correctness, verdict quality, and evidence compactness. This ordering
ensures that GRPO receives policy-gradient signals aligned with the desired
learning hierarchy.

\noindent\textbf{Constraint Derivation.}
We first derive constraints directly from the six design goals.

\noindent\emph{C1
(Format gate).} The worst format-valid trajectory has $o=0$ and
$V=-T_{\max}=-3$, so Goal 1 requires it to beat the format-failure penalty:
\[
-3w_v>c.
\]
\emph{C2 (Outcome dominance).} The worst correct trajectory must beat the
best wrong trajectory:
\[
R^{\inf}_{o=1}>R^{\sup}_{o=0}
\quad\Longrightarrow\quad
-3w_v+w_o>3w_v
\quad\Longrightarrow\quad
w_o>6w_v.
\]
\emph{C3 (Verdict discrimination and multi-round credit).} Since the reward
is linear in $V$ with slope $w_v$, verdict ordering and cumulative multi-round
credit require
\[
w_v>0.
\]
\emph{C4 (Compactness preference).} Since compactness only affects correct
trajectories through $w_l\ell$, compactness preference requires
\[
w_l>0.
\]
\emph{C5 (Verdict over length).} A single verdict correction changes
$V$ by $2$, so it must dominate the entire compactness range:
\[
2w_v>w_l.
\]

\noindent\textbf{Gap-Ordering Constraints.}
We next derive constraints from the desired hierarchy
$\Delta_1>\Delta_2>\Delta_3>\Delta_4>0$.

\noindent\emph{G1} follows from
$\Delta_1>\Delta_2$:
\[
-3w_v-c>w_o-6w_v
\quad\Longrightarrow\quad
c<3w_v-w_o.
\]
\emph{G2} follows from $\Delta_2>\Delta_3$:
\[
w_o-6w_v>2w_v
\quad\Longrightarrow\quad
w_o>8w_v.
\]
\emph{G3} follows from $\Delta_3>\Delta_4$:
\[
2w_v>w_l.
\]
\emph{G4} follows from $\Delta_4>0$:
\[
w_l>0.
\]

Combining the six ordering goals with the gap hierarchy gives the feasible
coefficient region:
\[
\mathcal{W}=
\left\{(c,w_v,w_o,w_l)\in\mathbb{R}\times\mathbb{R}_{>0}^{3}\;\middle|\;
\begin{aligned}
&w_v>0,\\
&w_o>8w_v,\\
&c<3w_v-w_o,\\
&2w_v>w_l
\end{aligned}
\right\}.
\]
This region keeps outcome correctness as the dominant task-level signal,
preserves verdict correctness as process-level supervision, and treats format
and length terms as stabilizing constraints rather than primary task objectives.

\noindent\textbf{Coefficient Selection.}
We select coefficients in three steps. First, we instantiate the design goals
as a four-level priority structure: answer correctness is the primary objective,
verdict correctness provides process-level control over stopping and continued
search, compactness improves efficiency, and format validity acts as a hard
gate. Second, this priority induces ratio constraints among weights: $w_o$
should dominate $w_v$, while the single-step verdict signal $2w_v$ should cover
the full length range $w_l$, preventing the policy from sacrificing verdict
quality for shorter evidence. Third, we choose the default setting
$(w_o,w_v,w_l)=(0.50,0.10,0.05)$. This makes outcome the dominant signal in our
rollout distribution, gives a verdict-step margin $2w_v=0.20$ larger than the
length range $w_l=0.05$, and sets the format penalty by
$c=3w_v-w_o=-0.20$ so that format failure remains below the worst valid output.

\noindent\textbf{Effectiveness Analysis.}
Our construction can be viewed as an expressibility argument for a desired
trajectory ordering. Following the perspective of
Abel et al.~\citep{abel2021expressivity}, reward design is not merely assigning
heuristic scalar scores: it asks whether a reward function can express a
task-level preference relation. In our case, the six design goals define a
partial order over rollout trajectories for temporal-causal evidence synthesis.
For example, outcome dominance requires every correct-answer trajectory to
outrank every wrong-answer trajectory, while the verdict and length goals impose
finer preferences within trajectory subsets. The feasible region $\mathcal{W}$
then gives a constructive certificate that this ordering is expressible by our
multi-level reward family. Moreover, the gap hierarchy
$\Delta_1>\Delta_2>\Delta_3>\Delta_4>0$ strengthens strict preference by
requiring larger guaranteed margins for higher-priority distinctions, echoing
the range-based view that acceptable behaviors should be separated from
unacceptable ones by reward margins. This explains why the coefficient design is
effective: it aligns the scalar GRPO learning signal with the intended
trajectory-level preference hierarchy, rather than relying on unstructured
reward mixing.
